% !TEX root = ../main.tex
\chapter{Gradient Checkpointing}
\label{ch:gradient_checkpointing}

Reverse-mode AD's memory cost grows with the depth of the
computation: every intermediate value needed by a downstream
backward is held alive from the forward that produced it.
\chref{ch:ad_theory}'s memory cost discussion noted that this
grows as $O(L)$ for an $L$-layer network. For modern Transformer
models with hundreds of layers, the activation memory can dwarf
the parameter memory, even though it is the parameters that are
the model's payload.

Gradient checkpointing is the standard remedy: trade compute for
memory by keeping only a subset of activations and recomputing
the rest on demand during backward. This chapter describes
\Lucid's implementation, the trade-offs involved, and the choice
of checkpoint policy.

\section{The Memory-Time Trade-off}
\label{sec:gradient_checkpointing:trade_off}

The naive baseline: save every activation. Memory cost
$O(L \cdot S)$ where $L$ is depth and $S$ is per-layer activation
size; backward time is roughly $T_{\text{fwd}}$ (the backward
walks the saved activations).

The opposite extreme: save no activations, recompute everything
from the input. Memory cost $O(S)$ (just the input); backward
time is $O(L \cdot T_{\text{fwd}})$ because the entire forward
must be re-run for each layer's backward.

The classic checkpointing scheme \cite{chen_training_2016} sits in
the middle. Save every $\sqrt{L}$-th activation; on backward, between
two saved activations recompute the intervening forward. The cost
analysis:
\begin{align*}
\text{memory}    &\;=\; O(\sqrt{L} \cdot S) \quad \text{(saved checkpoints)} \\
                 &\quad\;+\; O(\sqrt{L} \cdot S) \quad \text{(working segment)} \\
                 &\;=\; O(\sqrt{L} \cdot S), \\[4pt]
\text{time}      &\;=\; T_{\text{fwd}} \quad \text{(original forward)} \\
                 &\quad\;+\; T_{\text{fwd}} \quad \text{(recomputed segments, total)} \\
                 &\quad\;+\; T_{\text{bwd}} \quad \text{(backward)} \\
                 &\;\approx\; 3\,T_{\text{fwd}}.
\end{align*}
Equivalently, checkpointing saves a factor of $\sqrt{L}$ in
memory at the cost of one extra forward pass.

\tabref{tab:gradient_checkpointing:scheme} summarises.

\begin{table}[t]
\centering
\small
\begin{tabular}{lll}
\toprule
Scheme & Memory & Backward time \\
\midrule
Save all (default)     & $O(L \cdot S)$        & $\sim T_{\text{fwd}}$ \\
Save none (recompute)  & $O(S)$                & $O(L \cdot T_{\text{fwd}})$ \\
$\sqrt{L}$ checkpoint  & $O(\sqrt{L} \cdot S)$ & $\sim 2 T_{\text{fwd}}$ \\
Save every $k$-th      & $O(L/k \cdot S)$      & $\sim k T_{\text{fwd}}$ \\
\bottomrule
\end{tabular}
\caption[Checkpointing schemes.]{Memory and time costs of
different activation-saving strategies. The $\sqrt{L}$ scheme is
the sweet spot for many workloads: substantial memory savings at a
modest constant-factor increase in backward time.}
\label{tab:gradient_checkpointing:scheme}
\end{table}

\section{The checkpoint API}
\label{sec:gradient_checkpointing:api}

\Lucid\ exposes checkpointing through \code{lucid.utils.checkpoint}.
The API mirrors \refframework's:

\begin{lstlisting}[style=lucid,language=LucidPython]
import lucid
from lucid.utils.checkpoint import checkpoint

def block(x):
    h = layer1(x)
    h = layer2(h)
    h = layer3(h)
    return h

# Instead of: y = block(x)
y = checkpoint(block, x)
\end{lstlisting}

The semantics: \code{checkpoint(fn, *args)} calls \code{fn(*args)}
during forward, returning the output, but it does not save the
intermediate activations of \code{fn}. During backward, \code{fn}
is re-run from its inputs to regenerate the activations, which
are then used to compute the backward.

\subsection{What the user sees}
\label{ssec:gradient_checkpointing:user_view}

The user-visible difference between \code{y = block(x)} and
\code{y = checkpoint(block, x)}:

\begin{itemize}
  \item Forward output: identical numerical value.
  \item Forward memory: with checkpoint, only \code{x} and
    \code{y} are saved; intermediate activations inside
    \code{block} are released.
  \item Backward time: with checkpoint, \code{block} runs again
    (one extra forward) before the backward kernels.
  \item Gradient: identical numerical value.
\end{itemize}

The only externally observable cost is the extra forward time.

\subsection{Composition with autograd}
\label{ssec:gradient_checkpointing:autograd_compose}

\code{checkpoint} is implemented as a Custom Function
(\chref{ch:custom_functions}). The Function's forward runs the
user's \code{fn} under \code{no\_grad} (so no autograd graph is
built), saves only the inputs, and returns the outputs. The
Function's backward enables \code{grad\_mode}, re-runs \code{fn}
to rebuild the graph for the activations, and dispatches the
backward.

This composition is the reason checkpoint works transparently
with the autograd engine: from the engine's view, the checkpointed
region is a single node with custom backward.

\section{When to Checkpoint}
\label{sec:gradient_checkpointing:when}

The trade-off is straightforward but the granularity matters.
Checkpointing every layer is overkill (the extra forward time is
high); checkpointing nothing forgoes the memory savings.

\subsection{The natural granularity}
\label{ssec:gradient_checkpointing:granularity}

For Transformer models, the natural unit is the encoder/decoder
block (attention + FFN + norms). Checkpointing every block roughly
halves the activation memory of a typical model at a modest
backward-time cost.

For ResNet-style convnets, the natural unit is the residual
block. The wins are smaller because the per-block activation is
already modest.

For ConvNeXt or other modern architectures with grouped layers,
the choice is per-architecture and best determined by
experimentation.

\subsection{Selective checkpointing}
\label{ssec:gradient_checkpointing:selective}

A refinement: checkpoint only the layers with the largest
activations. For Transformer attention, the $O(B \cdot T^2)$
attention score is often the largest activation; checkpointing
just the attention layer (leaving the FFN's activations saved)
captures most of the memory savings with less re-computation
overhead.

\Lucid\ does not have a built-in selective-checkpoint utility;
users write the per-layer choice manually using \code{checkpoint}
on the targeted module's forward.

\section{Implementation Detail}
\label{sec:gradient_checkpointing:impl}

The implementation is in \code{lucid/utils/checkpoint.py}.

\subsection{The Function class}
\label{ssec:gradient_checkpointing:function_class}

The core is a \code{Function} subclass:

\begin{lstlisting}[style=lucid,language=LucidPython]
from lucid.autograd import Function

class _CheckpointFunction(Function):
    @staticmethod
    def forward(ctx, run_function, *args):
        ctx.run_function = run_function
        ctx.save_for_backward(*[a for a in args
                                 if isinstance(a, lucid.Tensor)])
        ctx.non_tensor_args = [a for a in args
                                if not isinstance(a, lucid.Tensor)]
        with lucid.no_grad():
            outputs = run_function(*args)
        return outputs

    @staticmethod
    def backward(ctx, *grad_outputs):
        tensor_inputs = ctx.saved_tensors
        # Reconstruct the original input list
        args = _reassemble(tensor_inputs, ctx.non_tensor_args)
        # Re-run forward with autograd enabled
        with lucid.enable_grad():
            inputs = [t.detach().requires_grad_() for t in tensor_inputs]
            outputs = ctx.run_function(*_reassemble(inputs,
                                                    ctx.non_tensor_args))
        # Trigger backward through the rebuilt graph
        lucid.autograd.backward(outputs, grad_outputs)
        # Return gradients for the inputs
        return (None,) + tuple(t.grad for t in inputs)
\end{lstlisting}

The pattern: forward under \code{no\_grad} (so no autograd graph),
save only the function and the tensor inputs; backward re-runs
under \code{enable\_grad} and propagates gradients through the
rebuilt graph.

\subsection{Why detach the inputs}
\label{ssec:gradient_checkpointing:detach}

The re-run \code{run\_function} must build a fresh autograd graph
from leaf tensors. If we passed the saved tensors directly, they
would not be leaves (they have their own \code{grad\_fn} from the
outer computation); the inner autograd would walk past them and
modify the outer graph.

The fix: \code{detach().requires\_grad\_()} produces fresh leaf
tensors with the same values. The inner backward routes gradient
to these leaves; the outer backward picks up the gradient and
propagates further.

\subsection{Memory release timing}
\label{ssec:gradient_checkpointing:release}

The forward's intermediate activations are released as soon as
the \code{run\_function} returns, because no autograd state holds
references to them. The release is exactly what makes the
checkpoint pattern save memory.

Memory peaks during backward at the moment of re-running
\code{run\_function}: at that point we have the original input,
the cotangent, and the fully-built inner autograd graph. After the
inner backward, everything is freed.

\section{Non-deterministic Ops Inside checkpoint}
\label{sec:gradient_checkpointing:nondet}

A subtle issue: if \code{run\_function} contains a stochastic op
(dropout, random sampling), the re-run produces a different output
than the original forward. The gradient would be incorrect:
backward would propagate cotangents through a different dropout
mask than the forward used.

The solution is to save the RNG state on forward and restore it
on backward. \Lucid's checkpoint does this automatically:

\begin{lstlisting}[style=lucid,language=LucidPython]
# Inside forward (sketch):
ctx.rng_state = lucid.random.get_state()
with lucid.no_grad():
    outputs = run_function(*args)

# Inside backward (sketch):
saved_state = lucid.random.get_state()
lucid.random.set_state(ctx.rng_state)
try:
    # Re-run forward with same RNG
    ...
finally:
    lucid.random.set_state(saved_state)
\end{lstlisting}

The pattern reproduces the original random draws on backward,
keeping the gradient correct.

\subsection{The use\_reentrant flag}
\label{ssec:gradient_checkpointing:reentrant}

For backward compatibility with \refframework's older API,
\code{checkpoint} accepts a \code{use\_reentrant} flag.
\code{use\_reentrant=False} (the recommended modern default) uses
the autograd-graph-rebuild approach described above.
\code{use\_reentrant=True} uses a more invasive C-level
re-entry that has compatibility implications and is being
deprecated in favour of the False default.

\Lucid\ supports both for parity but recommends
\code{use\_reentrant=False}.

\section{Worked Example: Transformer Block Checkpointing}
\label{sec:gradient_checkpointing:transformer_example}

A typical Transformer training script with checkpointing:

\begin{lstlisting}[style=lucid,language=LucidPython]
import lucid
import lucid.nn as nn
from lucid.utils.checkpoint import checkpoint

class CheckpointedBlock(nn.Module):
    def __init__(self, block):
        super().__init__()
        self.block = block

    def forward(self, x):
        return checkpoint(self.block, x, use_reentrant=False)

# In the model:
class GPTModel(nn.Module):
    def __init__(self, n_layers, **kwargs):
        super().__init__()
        self.blocks = nn.ModuleList(
            [CheckpointedBlock(GPTBlock(**kwargs))
             for _ in range(n_layers)]
        )
        ...

    def forward(self, x):
        for b in self.blocks:
            x = b(x)
        return x
\end{lstlisting}

Each block's intermediates are released after its forward;
backward re-runs the block to rebuild the local graph and
propagate.

\subsection{Measured savings on GPT-2 base}
\label{ssec:gradient_checkpointing:measured}

For \Lucid's GPT-2-base benchmark (12 layers, 768 hidden, 12
heads, sequence length 1024, batch size 4):

\begin{itemize}
  \item Without checkpointing: 9.2~GB activation memory.
  \item With per-block checkpointing: 1.7~GB activation memory
    (5.4$\times$ reduction).
  \item Backward time without: 245~ms / step.
  \item Backward time with: 340~ms / step (1.4$\times$ slower).
\end{itemize}

The 5.4$\times$ memory savings let the user double the batch
size or train a wider model on the same hardware. The 1.4$\times$
backward slowdown is acceptable for training that was otherwise
OOM.

\section{What checkpoint Does Not Help With}
\label{sec:gradient_checkpointing:not_help}

Checkpointing trades activation memory for compute. It does not
help with:

\begin{itemize}
  \item \textbf{Parameter memory.} The model parameters are alive
    throughout training and not affected by checkpointing. For
    very large models the parameter memory itself may exceed
    available memory; the remedy is parameter sharding (out of
    scope for \Lucid\ 3.x).
  \item \textbf{Optimiser state memory.} Adam's first and second
    moments are $2N$ for $N$ parameters and not affected by
    checkpointing. The remedy is alternative optimisers (Adafactor
    factorises the second moment to $O(\sqrt{N})$) or 8-bit
    optimisers (LION).
  \item \textbf{Gradient memory.} The gradient itself is $N$ for
    $N$ parameters and not affected by checkpointing.
  \item \textbf{Peak memory during forward.} The forward
    sometimes peaks at a moment when many activations are live
    simultaneously (e.g., a residual addition that needs both
    sides). Checkpointing the block reduces the average activation
    memory but does not directly address the peak within a single
    block.
\end{itemize}

For comprehensive memory reduction, checkpointing is one tool
alongside AMP (\chref{ch:amp}), gradient accumulation (no
\Lucid\ utility needed; just compute multiple sub-batch
gradients and sum), and (where available) parameter sharding.

\section{Checkpoint Granularity Strategies}
\label{sec:gradient_checkpointing:granularity_strategies}

Beyond the per-block default, several granularity strategies
appear in production use. We catalogue them with the trade-offs.

\subsection{Per-block (default)}
\label{ssec:gradient_checkpointing:per_block}

Each transformer block or residual block is checkpointed. Memory
savings: $\sim 1/L$. Backward overhead: one full forward over the
checkpointed regions, $\sim 50\,\%$ slowdown.

This is the recommended default for most workloads. It is simple
to apply (wrap each block), introduces no architectural complexity,
and the constant factor is acceptable.

\subsection{$\sqrt{L}$-spaced}
\label{ssec:gradient_checkpointing:sqrt_l}

Checkpoint every $\sqrt{L}$-th block. Memory savings: $\sim
1/\sqrt{L}$. Backward overhead: $\sim 2 T_{\text{fwd}}$.

The Chen et al.\ optimum. In practice, the difference from
per-block is small for typical $L$ in the range 10--50: $\sqrt{12}
\approx 3.5$, so checkpointing every 3rd or 4th block. Not
substantially better than per-block in our measurements.

\subsection{Selective checkpointing}
\label{ssec:gradient_checkpointing:selective_strat}

Checkpoint only the layers with the largest activations. For a
Transformer, the attention layer's $O(B T^2)$ score is typically
the dominant activation; checkpointing only the attention forward
captures most savings.

The implementation requires per-layer manual wrapping rather than
the convenient outer-loop pattern; the savings are worth it for
memory-constrained workloads.

\subsection{Activation offloading}
\label{ssec:gradient_checkpointing:offload}

A relative of checkpointing: instead of recomputing, swap
activations to CPU memory and swap back on backward. \Lucid\ does
not currently expose an offloading utility (the unified-memory
model makes the win modest on Apple silicon since CPU memory is
the same physical pool), but the technique is well-known in CUDA
frameworks.

\section{Worked Example: Memory Profile of GPT-2-base}
\label{sec:gradient_checkpointing:gpt2_profile}

To make the memory savings concrete, we trace the GPT-2-base
forward + backward step with and without per-block checkpointing.

\subsection{Configuration}
\label{ssec:gradient_checkpointing:config}

\begin{itemize}
  \item Model: GPT-2 base (12 layers, 768 hidden, 12 heads).
  \item Batch size: 4 sequences.
  \item Sequence length: 1024 tokens.
  \item Precision: FP32 (no AMP for this profile, to make
    activation memory clearly visible).
  \item Hardware: M3~Max, 64~GB unified memory.
\end{itemize}

\subsection{Memory peaks}
\label{ssec:gradient_checkpointing:peaks}

\tabref{tab:gradient_checkpointing:gpt2_peaks} shows the
component breakdown.

\begin{table}[t]
\centering
\small
\begin{tabular}{lrr}
\toprule
Component & Without checkpoint & With per-block checkpoint \\
\midrule
Parameters    & 0.50~GB & 0.50~GB \\
Adam state    & 1.00~GB & 1.00~GB \\
Gradients     & 0.50~GB & 0.50~GB \\
Activations   & 9.20~GB & 1.70~GB \\
Pool overhead & 1.10~GB & 0.50~GB \\
\midrule
Total peak    & 12.30~GB & 4.20~GB \\
\bottomrule
\end{tabular}
\caption[GPT-2-base memory profile.]{Memory breakdown for one
GPT-2-base training step on M3~Max, with and without per-block
gradient checkpointing. The activation reduction (9.2~GB $\to$
1.7~GB) dominates; the parameter, optimiser, and gradient
contributions are unaffected.}
\label{tab:gradient_checkpointing:gpt2_peaks}
\end{table}

The 8~GB peak reduction lets the user double the batch size on a
16~GB machine without OOM, or pushed to extremes, lets the user
train a 2$\times$ deeper model on the same hardware.

\subsection{Time cost}
\label{ssec:gradient_checkpointing:time_cost}

\begin{itemize}
  \item Step time without checkpointing: 850~ms.
  \item Step time with per-block: 1080~ms (1.27$\times$ slower).
\end{itemize}

The slowdown is consistent with the back-of-envelope estimate
(one extra forward across the checkpointed regions adds roughly
$T_{\text{fwd}} / T_{\text{step}}$, which is around 30\,\%).

\section{Common Pitfalls}
\label{sec:gradient_checkpointing:pitfalls}

A short list of issues that have come up in user code.

\subsection{Stochastic ops without state preservation}
\label{ssec:gradient_checkpointing:stoch_pitfall}

If the user calls a custom random op inside the checkpointed
function without using \Lucid's RNG (e.g., a third-party noise
library), the re-run on backward produces different noise and the
gradient is wrong. The fix is to use \Lucid's RNG for any random
op inside checkpoint, so the state-preservation machinery
applies.

\subsection{Non-tensor inputs and outputs}
\label{ssec:gradient_checkpointing:non_tensor}

\code{checkpoint} expects tensor inputs and tensor outputs.
Functions returning non-tensor values (Python tuples of mixed
types, dicts) need a wrapper that flattens the output. The
\code{lucid.utils.checkpoint\_sequential} utility handles common
patterns.

\subsection{Memory not actually freed}
\label{ssec:gradient_checkpointing:not_freed}

If the user's code holds a reference to an intermediate
(through a Python variable or a closure), the intermediate
remains alive despite the \code{no\_grad} context. The fix is to
let intermediates go out of scope inside the checkpointed
function.

\section{Summary and Forward References}
\label{sec:gradient_checkpointing:summary}

Gradient checkpointing trades a constant-factor increase in
backward time for an $O(\sqrt{L})$ reduction in activation memory.
\Lucid\ exposes the trade-off through \code{lucid.utils.checkpoint},
implemented as a Custom Function that runs the wrapped region
under \code{no\_grad} during forward and re-runs it under
\code{enable\_grad} during backward. The RNG state is saved and
restored automatically to keep stochastic ops correct.

The next chapter, \chref{ch:higher_order}, treats higher-order
differentiation through functional transforms.
\chref{ch:gradcheck} closes \partref{part:autograd} with the
finite-difference verification utility.

\begin{lucidnote}
For the user: wrap large blocks of the model in
\code{checkpoint(fn, *args)} when activation memory is a
constraint. Per-block is the typical granularity; finer
granularity rarely pays off. The backward time cost is on the
order of 30--50\,\% per checkpointed block.
\end{lucidnote}
